\section{Related Work}

Variational Autoencoders (VAEs) have been extensively used for representation learning \cite{kingma2019introduction} as their amortized inference network lends itself naturally to this task.
In the context of automated planning, \citet{DBLP:journals/corr/AsaiF17} proposed a State Autoencoder (SAE) for propositional symbolic grounding. SAEs are in fact VAEs with Bernoulli latent variables. They are trained by maximizing a modified ELBO with an additional entropy regularizer, defined as twice the negative KL divergence. Thus, the objective function being maximized is the ELBO~\eqref{eq:elbo} with the sign of the KL flipped. 
%Although unintentional \cite{asai2019stable}, this proved to be fundamental for the mitigation of the issue of representation instability.
%
A variation of SAE, the zero-suppressed state autoencoder \cite{asai2019stable}, adds a further regularization term to the propositional representation, leading to more stable features.

\citet{zhang2018composable} learn state features with supervised learning, and then plan in the feature space with a learned transition graph. 
\citet{konidaris2018skills} specify a set of skills for a task, and automatically extract state representations from raw observations. 
\citet{kurutach2018learning} use generative adversarial networks to learn structured representations of images and a deterministic dynamics model, and plan with graph search.
%
\citet{RLIW} proposed $\pi$-IW, a variant of RolloutIW where a neural network guides the action selection process in the rollouts, which would otherwise be random.
Moreover, $\pi$-IW plans using features obtained from the last hidden layer of the policy network, instead of B-PROST.

Atari games have also been extensively used as a testbed in model-free \citep{dqn, dqn2,a3c,ppo,ga3c,acktr,vtrace, rainbow,kapturowski2018recurrent} and model-based \citep{vpn,dyna_dqn,kaiser2019model,schrittwieser2019mastering,badia2020agent57} deep reinforcement learning. 